\chapter*{来源对应与复现}
\addcontentsline{toc}{chapter}{来源对应与复现}
本稿对五组本地理论来源按数学依赖重新组织。外部散度、PMD-Mean 和 TOP-D 资料在对应位置引用；本地来源保留原文件，以便追溯定义及历史语境。同一命题在多个版本中出现时，本稿只维护一个完整证明。

随源文件提供的 \path{source-map.csv} 记录原文件、命题位置、源文件校验和及新稿主题归属。主题合并可能是多对一关系；对应表用于追溯材料，不将重复版本计为独立理论结果。原文没有命题环境的推导另以文件条目登记。

\begin{longtable}{p{24mm}p{57mm}p{47mm}}
\toprule
来源目录 & 主要材料 & 新稿归属\\
\midrule
\endhead
\path{alpha-divergence/} & 密度比矩、散度端点、score 梯度、几何重加权测度 & 附录 A、B、D；统一教师记号与 Hellinger 常数\\
\path{book/} & 重心、$J$、四种 inner、信息访问、估计、动力学与边界 & 正文及附录 A--H 的主要数学来源\\
\path{book-old/} & Mean 修订、KL 控制族、几何与 Mean 的同步推导 & 附录 C、E、F、G；保留历史差异的数学内容\\
\path{Notes/} & 整合证明、成对 BPO、无重复动作、方差与误差连接 & 附录 B、D、F、G，特别是 U-统计量完整分解\\
\path{paper/} & 参数化叙事、曲率、受限投影与有限步论证 & 正文结构及附录 C、D；精确展开曲率与驻点证明\\
\bottomrule
\end{longtable}

来源中的纯实验配置、旧生产模块路径与任务成绩保留在原材料中。本稿的反例和验证针对所写数学对象。来源名相同不表示统计口径相同：头部可能由教师或锚点选出，条件动作采样可能来自当前学生或旧锚点，必须按本稿每条定理的条件代入。

一些跨主题内容的归属为：多教师重心与通用单调信号在附录 B；独立拟合形状、共享参数投影和标量谱调度在 C、G；最坏矩、相对噪声和绝对噪声准则在 D；配分函数、自动微分与非保守场在 E、F；教师漂移、同步及序列反例在 G。幂散度梯度中的几何重加权分布 $p(T/p)^\kappa$ 与本轮幂均值目标 $m_{\kappa,\eta}$ 是不同对象，前者用于积分测度解释，后者由重心问题定义。

\section*{数学核验的范围}
数值核验采用有限词表的确定性求和，对照解析恒等式、自动微分和独立的概率公式。重点包括重心超额损失、$J$ 的梯度与 JS 端点、四种 inner 的梯度和 Hessian、成对估计的协方差、形状端点、TOP-D 算术收缩以及条件与序列的反例。测试入口为 \path{tests/reference/test_core_monograph.py}。

2026 年 10 月 15 日的 CPU 核验中，该入口的 46 个参数化测试全部通过；仓库全套测试为 316 项通过。积分目标采用 96 点 Gauss--Legendre 求积与解析梯度对照，主要恒等式容差为 $10^{-11}$ 或 $10^{-12}$。U-统计量使用三动作词表、$M=2,3,4$ 的全部有序样本组合；无重复抽样直接比较协方差矩阵，检查半正定差值。局部三次率采用逐渐接近教师的固定方向，不以有限差分噪声替代极限系数。运行环境与命令保存于 \path{validation.json}。

数值测试用于暴露符号、尺度和实现错误；有限实例通过不构成定理证明。理论正文中的结论以相应解析证明和显式条件为依据。测试结果不用于声称共享参数大模型的训练收敛或任务准确率。


在仓库根执行 \texttt{make theory-core} 构建完整稿和正文阅读版；在本目录执行 \texttt{make} 等价。构建复用相邻 \path{../book/elegantbook.cls}，使用 XeLaTeX、latexmk 与 Biber。正文阅读版由同一主源通过条件装配生成，附录引用保留完整稿的编号。
